% =====================================================================
% Paper 10 — Pillar-10 canonical structural hbar formula + GAP-1
%            verification programme paper
%   (HEALTHIEST of Wave 1/4/5 batch; canonical Math291 formula preserved;
%    light wording calibration only)
% Status: [DRAFT] [NEEDS_UPDATE]
% AUDIT-FLAG: AUDIT-2026-05-02-Wave7-Aux-Epoch-Overclaim
%   (Math314-AddD extension covering Wave 1/4/5 papers Paper-09..16).
%   Hostile-referee audit (2026-05-02): the BEST-positioned paper of
%   the Wave 1/4/5 batch. Canonical Math291 formula preserved verbatim;
%   classical-no-go-theorem-then-KZ-resolution structure aligned with
%   GAP-1 programme. Light wording correction only:
%   "matter-side substantiation confirms the absence of competing
%   contributions above 10^-3" → conditional/programme framing per
%   Math297 / Math299 / Math312 verdict-shell pending.
% Compile: pdflatex Paper-10.tex (×2 for refs)
% Canonical sources: Math79-AddB, Math98-AddA..E, Math110-AddI, Math163,
%                    Math191, Math196, Math291, Math297, Math299,
%                    Math312, Math314-AddD
% =====================================================================
% --- TECT canonical preamble (see _shared/tect-preamble.tex) ----------
% =====================================================================
% TECT canonical LaTeX preamble (revtex4-2 / single-column / theorem-safe)
% =====================================================================
% Maintainer: Jusang Lee (jtkor@outlook.com)
% Binding from: 2026-05-06
% Purpose: a single source-of-truth preamble for every TECT paper
% (Paper-00 .. Paper-10 + Auxiliary notes). Designed to (a) eliminate
% font-corruption on pdflatex (T1 fontenc + lmodern), (b) make
% theorem-heavy mathematical-physics content render cleanly in a
% single-column layout (PRD class option, NOT PRL two-column), (c)
% provide consistent theorem numbering across all TECT papers via
% amsthm with a shared counter set, (d) make hyperref + cleveref
% cross-references resolve cleanly with the standard two-pass
% pdflatex compile.
%
% Usage in each Paper-NN.tex:
%   % =====================================================================
% TECT canonical LaTeX preamble (revtex4-2 / single-column / theorem-safe)
% =====================================================================
% Maintainer: Jusang Lee (jtkor@outlook.com)
% Binding from: 2026-05-06
% Purpose: a single source-of-truth preamble for every TECT paper
% (Paper-00 .. Paper-10 + Auxiliary notes). Designed to (a) eliminate
% font-corruption on pdflatex (T1 fontenc + lmodern), (b) make
% theorem-heavy mathematical-physics content render cleanly in a
% single-column layout (PRD class option, NOT PRL two-column), (c)
% provide consistent theorem numbering across all TECT papers via
% amsthm with a shared counter set, (d) make hyperref + cleveref
% cross-references resolve cleanly with the standard two-pass
% pdflatex compile.
%
% Usage in each Paper-NN.tex:
%   % =====================================================================
% TECT canonical LaTeX preamble (revtex4-2 / single-column / theorem-safe)
% =====================================================================
% Maintainer: Jusang Lee (jtkor@outlook.com)
% Binding from: 2026-05-06
% Purpose: a single source-of-truth preamble for every TECT paper
% (Paper-00 .. Paper-10 + Auxiliary notes). Designed to (a) eliminate
% font-corruption on pdflatex (T1 fontenc + lmodern), (b) make
% theorem-heavy mathematical-physics content render cleanly in a
% single-column layout (PRD class option, NOT PRL two-column), (c)
% provide consistent theorem numbering across all TECT papers via
% amsthm with a shared counter set, (d) make hyperref + cleveref
% cross-references resolve cleanly with the standard two-pass
% pdflatex compile.
%
% Usage in each Paper-NN.tex:
%   \input{../_shared/tect-preamble.tex}
%   \begin{document}
%   ...
%   \end{document}
%
% Build (every paper): `pdflatex Paper-NN && pdflatex Paper-NN`
% (the second pass resolves \ref{} / \cref{} forward references).
% A bash/PowerShell wrapper that automates this is at
% Docs/papers/papers/_shared/build-paper.sh / build-paper.ps1.
% =====================================================================

% ---- Document class --------------------------------------------------
% PRD class option of revtex4-2 with [reprint,onecolumn] gives the same
% APS heading / by-line / abstract style as PRL but in a wider single
% column that handles long display equations and theorem environments
% without column-break artefacts. nofootinbib keeps citations out of
% footnotes. The 11pt option restores standard font metrics; with T1
% fontenc + lmodern this gives non-bitmap glyphs.
\documentclass[%
  aps,prd,reprint,onecolumn,nofootinbib,11pt,%
  amsmath,amssymb,floatfix,longbibliography%
]{revtex4-2}

% ---- Encoding & fonts ------------------------------------------------
\usepackage[T1]{fontenc}        % proper 8-bit font encoding (no font corruption)
\usepackage[utf8]{inputenc}     % allow UTF-8 source
\usepackage{lmodern}            % scalable Latin Modern (replaces bitmap CM)
\usepackage{textcomp}           % extra text symbols
\usepackage{microtype}          % subtle kerning improvements

% ---- UTF-8 → LaTeX character mappings --------------------------------
% Maps the Unicode characters that occur in TECT body text (legacy from
% the chat-content auto-archival rule, CLAUDE.md §4) to LaTeX-safe
% commands. Without these declarations, T1+inputenc rejects the chars
% with "Unicode character X not set up for use with LaTeX".
% Adding declarations centrally (here) is preferable to per-file edits.
\DeclareUnicodeCharacter{00A7}{\S}              % §  (section sign)
\DeclareUnicodeCharacter{00B2}{\ensuremath{^{2}}}        % ²
\DeclareUnicodeCharacter{00B3}{\ensuremath{^{3}}}        % ³
\DeclareUnicodeCharacter{00B5}{\ensuremath{\mu}}         % µ (micro)
\DeclareUnicodeCharacter{00B7}{\ensuremath{\cdot}}       % · (middle dot)
\DeclareUnicodeCharacter{00D7}{\ensuremath{\times}}      % ×
\DeclareUnicodeCharacter{00E4}{\"a}             % ä
\DeclareUnicodeCharacter{00E9}{\'e}             % é
\DeclareUnicodeCharacter{00F6}{\"o}             % ö
\DeclareUnicodeCharacter{00FC}{\"u}             % ü
\DeclareUnicodeCharacter{010C}{\v{C}}           % Č
\DeclareUnicodeCharacter{010D}{\v{c}}           % č
\DeclareUnicodeCharacter{0394}{\ensuremath{\Delta}}      % Δ
\DeclareUnicodeCharacter{03A3}{\ensuremath{\Sigma}}      % Σ
\DeclareUnicodeCharacter{03A9}{\ensuremath{\Omega}}      % Ω
\DeclareUnicodeCharacter{03B1}{\ensuremath{\alpha}}      % α
\DeclareUnicodeCharacter{03B2}{\ensuremath{\beta}}       % β
\DeclareUnicodeCharacter{03B3}{\ensuremath{\gamma}}      % γ
\DeclareUnicodeCharacter{03B4}{\ensuremath{\delta}}      % δ
\DeclareUnicodeCharacter{03B5}{\ensuremath{\varepsilon}} % ε
\DeclareUnicodeCharacter{03BA}{\ensuremath{\kappa}}      % κ
\DeclareUnicodeCharacter{03BB}{\ensuremath{\lambda}}     % λ
\DeclareUnicodeCharacter{03BC}{\ensuremath{\mu}}         % μ
\DeclareUnicodeCharacter{03BD}{\ensuremath{\nu}}         % ν
\DeclareUnicodeCharacter{03C0}{\ensuremath{\pi}}         % π
\DeclareUnicodeCharacter{03C1}{\ensuremath{\rho}}        % ρ
\DeclareUnicodeCharacter{03C3}{\ensuremath{\sigma}}      % σ
\DeclareUnicodeCharacter{03C4}{\ensuremath{\tau}}        % τ
\DeclareUnicodeCharacter{03C6}{\ensuremath{\varphi}}     % φ
\DeclareUnicodeCharacter{03C7}{\ensuremath{\chi}}        % χ
\DeclareUnicodeCharacter{03C8}{\ensuremath{\psi}}        % ψ
\DeclareUnicodeCharacter{03C9}{\ensuremath{\omega}}      % ω
\DeclareUnicodeCharacter{2013}{--}              % – (en dash)
\DeclareUnicodeCharacter{2014}{---}             % — (em dash)
\DeclareUnicodeCharacter{2018}{`}               % ‘
\DeclareUnicodeCharacter{2019}{'}               % ’
\DeclareUnicodeCharacter{201C}{``}              % “
\DeclareUnicodeCharacter{201D}{''}              % ”
\DeclareUnicodeCharacter{2070}{\ensuremath{^{0}}}        % ⁰
\DeclareUnicodeCharacter{2071}{\ensuremath{^{i}}}        % ⁱ
\DeclareUnicodeCharacter{2122}{\textsuperscript{TM}}     % ™
\DeclareUnicodeCharacter{2126}{\ensuremath{\Omega}}      % Ω (ohm sign)
\DeclareUnicodeCharacter{210F}{\ensuremath{\hbar}}       % ℏ
\DeclareUnicodeCharacter{2115}{\ensuremath{\mathbb{N}}}  % ℕ
\DeclareUnicodeCharacter{2102}{\ensuremath{\mathbb{C}}}  % ℂ
\DeclareUnicodeCharacter{211D}{\ensuremath{\mathbb{R}}}  % ℝ
\DeclareUnicodeCharacter{2124}{\ensuremath{\mathbb{Z}}}  % ℤ
\DeclareUnicodeCharacter{211A}{\ensuremath{\mathbb{Q}}}  % ℚ
\DeclareUnicodeCharacter{2192}{\ensuremath{\to}}         % →
\DeclareUnicodeCharacter{21D2}{\ensuremath{\Rightarrow}} % ⇒
\DeclareUnicodeCharacter{2200}{\ensuremath{\forall}}     % ∀
\DeclareUnicodeCharacter{2203}{\ensuremath{\exists}}     % ∃
\DeclareUnicodeCharacter{2208}{\ensuremath{\in}}         % ∈
\DeclareUnicodeCharacter{2209}{\ensuremath{\notin}}      % ∉
\DeclareUnicodeCharacter{2227}{\ensuremath{\wedge}}      % ∧
\DeclareUnicodeCharacter{2228}{\ensuremath{\vee}}        % ∨
\DeclareUnicodeCharacter{2229}{\ensuremath{\cap}}        % ∩
\DeclareUnicodeCharacter{222A}{\ensuremath{\cup}}        % ∪
\DeclareUnicodeCharacter{2245}{\ensuremath{\cong}}       % ≅
\DeclareUnicodeCharacter{2248}{\ensuremath{\approx}}     % ≈
\DeclareUnicodeCharacter{2260}{\ensuremath{\neq}}        % ≠
\DeclareUnicodeCharacter{2264}{\ensuremath{\leq}}        % ≤
\DeclareUnicodeCharacter{2265}{\ensuremath{\geq}}        % ≥
\DeclareUnicodeCharacter{22C5}{\ensuremath{\cdot}}       % ⋅
\DeclareUnicodeCharacter{2295}{\ensuremath{\oplus}}      % ⊕
\DeclareUnicodeCharacter{2297}{\ensuremath{\otimes}}     % ⊗

% ---- Mathematics & symbols ------------------------------------------
\usepackage{amsmath,amssymb,bm,mathrsfs,mathtools}
\usepackage{amsthm}             % theorem environments (load BEFORE hyperref)

% ---- Theorem environments (shared counter set, TECT canonical) -------
% All theorems / lemmas / propositions / corollaries / remarks share a
% single counter so cross-references read consistently across a paper.
% Definitions get a separate counter (definitions are numbered in their
% own sequence by editorial convention).
\newtheorem{theorem}{Theorem}
\newtheorem{lemma}[theorem]{Lemma}
\newtheorem{proposition}[theorem]{Proposition}
\newtheorem{corollary}[theorem]{Corollary}

\theoremstyle{remark}
\newtheorem{remark}[theorem]{Remark}
\newtheorem{example}[theorem]{Example}

\theoremstyle{definition}
\newtheorem{definition}[theorem]{Definition}
\newtheorem{conjecture}[theorem]{Conjecture}
\newtheorem{hypothesis}[theorem]{Hypothesis}

% ---- Tables / floats / graphics --------------------------------------
\usepackage{booktabs}           % \toprule / \midrule / \bottomrule
\usepackage{array}
\usepackage{graphicx}
\usepackage{enumitem}           % \begin{itemize}[leftmargin=...] options
\usepackage{hyphenat}           % \hyp{} for breakable hyphens in \texttt
\usepackage{seqsplit}           % \seqsplit{} for long unbreakable strings
\sloppy                         % allow stretchier inter-word spacing to
                                % reduce overfull \hbox warnings on long URLs
                                % and long \texttt tokens (paragraph-level).
\emergencystretch=3em           % last-resort hbox stretch budget

% ---- Cross-references (hyperref last among the standard packages,
%      cleveref AFTER hyperref) ----------------------------------------
\usepackage[%
  colorlinks=true,%
  linkcolor=blue,%
  citecolor=blue,%
  urlcolor=blue,%
  breaklinks=true,%
  bookmarks=true,%
  bookmarksnumbered=true%
]{hyperref}
\usepackage[capitalise,nameinlink]{cleveref}

% ---- TECT-canonical author block macros -----------------------------
% (defined here so every paper has the same author / affiliation style)
\newcommand{\tectauthor}{%
  \author{Jusang Lee}%
  \email{jtkor@outlook.com}%
  \affiliation{Independent researcher, Republic of Korea}%
}

% ---- TECT-canonical archive footer ----------------------------------
\newcommand{\tectarchivefooter}{%
  \par\medskip
  \noindent\small\textit{Canonical archive}:
  \url{https://github.com/TECT-OpenPhysics/TECT}.\par%
}

% ---- TECT TODO / AUDIT-FLAG / HONEST-SCOPE inline marks --------------
% Used in drafts to highlight pending audit items without disturbing
% the body text style. Suppress via \tectsuppressmarks (define empty).
\providecommand{\tectauditflag}[1]{\textcolor{red}{\textbf{[AUDIT-FLAG: #1]}}}
\providecommand{\tecttodo}[1]{\textcolor{orange}{\textbf{[TODO: #1]}}}
\providecommand{\tecthonestscope}[1]{\textit{Honest scope: #1.}}

% =====================================================================
% End of TECT canonical preamble
% =====================================================================

%   \begin{document}
%   ...
%   \end{document}
%
% Build (every paper): `pdflatex Paper-NN && pdflatex Paper-NN`
% (the second pass resolves \ref{} / \cref{} forward references).
% A bash/PowerShell wrapper that automates this is at
% Docs/papers/papers/_shared/build-paper.sh / build-paper.ps1.
% =====================================================================

% ---- Document class --------------------------------------------------
% PRD class option of revtex4-2 with [reprint,onecolumn] gives the same
% APS heading / by-line / abstract style as PRL but in a wider single
% column that handles long display equations and theorem environments
% without column-break artefacts. nofootinbib keeps citations out of
% footnotes. The 11pt option restores standard font metrics; with T1
% fontenc + lmodern this gives non-bitmap glyphs.
\documentclass[%
  aps,prd,reprint,onecolumn,nofootinbib,11pt,%
  amsmath,amssymb,floatfix,longbibliography%
]{revtex4-2}

% ---- Encoding & fonts ------------------------------------------------
\usepackage[T1]{fontenc}        % proper 8-bit font encoding (no font corruption)
\usepackage[utf8]{inputenc}     % allow UTF-8 source
\usepackage{lmodern}            % scalable Latin Modern (replaces bitmap CM)
\usepackage{textcomp}           % extra text symbols
\usepackage{microtype}          % subtle kerning improvements

% ---- UTF-8 → LaTeX character mappings --------------------------------
% Maps the Unicode characters that occur in TECT body text (legacy from
% the chat-content auto-archival rule, CLAUDE.md §4) to LaTeX-safe
% commands. Without these declarations, T1+inputenc rejects the chars
% with "Unicode character X not set up for use with LaTeX".
% Adding declarations centrally (here) is preferable to per-file edits.
\DeclareUnicodeCharacter{00A7}{\S}              % §  (section sign)
\DeclareUnicodeCharacter{00B2}{\ensuremath{^{2}}}        % ²
\DeclareUnicodeCharacter{00B3}{\ensuremath{^{3}}}        % ³
\DeclareUnicodeCharacter{00B5}{\ensuremath{\mu}}         % µ (micro)
\DeclareUnicodeCharacter{00B7}{\ensuremath{\cdot}}       % · (middle dot)
\DeclareUnicodeCharacter{00D7}{\ensuremath{\times}}      % ×
\DeclareUnicodeCharacter{00E4}{\"a}             % ä
\DeclareUnicodeCharacter{00E9}{\'e}             % é
\DeclareUnicodeCharacter{00F6}{\"o}             % ö
\DeclareUnicodeCharacter{00FC}{\"u}             % ü
\DeclareUnicodeCharacter{010C}{\v{C}}           % Č
\DeclareUnicodeCharacter{010D}{\v{c}}           % č
\DeclareUnicodeCharacter{0394}{\ensuremath{\Delta}}      % Δ
\DeclareUnicodeCharacter{03A3}{\ensuremath{\Sigma}}      % Σ
\DeclareUnicodeCharacter{03A9}{\ensuremath{\Omega}}      % Ω
\DeclareUnicodeCharacter{03B1}{\ensuremath{\alpha}}      % α
\DeclareUnicodeCharacter{03B2}{\ensuremath{\beta}}       % β
\DeclareUnicodeCharacter{03B3}{\ensuremath{\gamma}}      % γ
\DeclareUnicodeCharacter{03B4}{\ensuremath{\delta}}      % δ
\DeclareUnicodeCharacter{03B5}{\ensuremath{\varepsilon}} % ε
\DeclareUnicodeCharacter{03BA}{\ensuremath{\kappa}}      % κ
\DeclareUnicodeCharacter{03BB}{\ensuremath{\lambda}}     % λ
\DeclareUnicodeCharacter{03BC}{\ensuremath{\mu}}         % μ
\DeclareUnicodeCharacter{03BD}{\ensuremath{\nu}}         % ν
\DeclareUnicodeCharacter{03C0}{\ensuremath{\pi}}         % π
\DeclareUnicodeCharacter{03C1}{\ensuremath{\rho}}        % ρ
\DeclareUnicodeCharacter{03C3}{\ensuremath{\sigma}}      % σ
\DeclareUnicodeCharacter{03C4}{\ensuremath{\tau}}        % τ
\DeclareUnicodeCharacter{03C6}{\ensuremath{\varphi}}     % φ
\DeclareUnicodeCharacter{03C7}{\ensuremath{\chi}}        % χ
\DeclareUnicodeCharacter{03C8}{\ensuremath{\psi}}        % ψ
\DeclareUnicodeCharacter{03C9}{\ensuremath{\omega}}      % ω
\DeclareUnicodeCharacter{2013}{--}              % – (en dash)
\DeclareUnicodeCharacter{2014}{---}             % — (em dash)
\DeclareUnicodeCharacter{2018}{`}               % ‘
\DeclareUnicodeCharacter{2019}{'}               % ’
\DeclareUnicodeCharacter{201C}{``}              % “
\DeclareUnicodeCharacter{201D}{''}              % ”
\DeclareUnicodeCharacter{2070}{\ensuremath{^{0}}}        % ⁰
\DeclareUnicodeCharacter{2071}{\ensuremath{^{i}}}        % ⁱ
\DeclareUnicodeCharacter{2122}{\textsuperscript{TM}}     % ™
\DeclareUnicodeCharacter{2126}{\ensuremath{\Omega}}      % Ω (ohm sign)
\DeclareUnicodeCharacter{210F}{\ensuremath{\hbar}}       % ℏ
\DeclareUnicodeCharacter{2115}{\ensuremath{\mathbb{N}}}  % ℕ
\DeclareUnicodeCharacter{2102}{\ensuremath{\mathbb{C}}}  % ℂ
\DeclareUnicodeCharacter{211D}{\ensuremath{\mathbb{R}}}  % ℝ
\DeclareUnicodeCharacter{2124}{\ensuremath{\mathbb{Z}}}  % ℤ
\DeclareUnicodeCharacter{211A}{\ensuremath{\mathbb{Q}}}  % ℚ
\DeclareUnicodeCharacter{2192}{\ensuremath{\to}}         % →
\DeclareUnicodeCharacter{21D2}{\ensuremath{\Rightarrow}} % ⇒
\DeclareUnicodeCharacter{2200}{\ensuremath{\forall}}     % ∀
\DeclareUnicodeCharacter{2203}{\ensuremath{\exists}}     % ∃
\DeclareUnicodeCharacter{2208}{\ensuremath{\in}}         % ∈
\DeclareUnicodeCharacter{2209}{\ensuremath{\notin}}      % ∉
\DeclareUnicodeCharacter{2227}{\ensuremath{\wedge}}      % ∧
\DeclareUnicodeCharacter{2228}{\ensuremath{\vee}}        % ∨
\DeclareUnicodeCharacter{2229}{\ensuremath{\cap}}        % ∩
\DeclareUnicodeCharacter{222A}{\ensuremath{\cup}}        % ∪
\DeclareUnicodeCharacter{2245}{\ensuremath{\cong}}       % ≅
\DeclareUnicodeCharacter{2248}{\ensuremath{\approx}}     % ≈
\DeclareUnicodeCharacter{2260}{\ensuremath{\neq}}        % ≠
\DeclareUnicodeCharacter{2264}{\ensuremath{\leq}}        % ≤
\DeclareUnicodeCharacter{2265}{\ensuremath{\geq}}        % ≥
\DeclareUnicodeCharacter{22C5}{\ensuremath{\cdot}}       % ⋅
\DeclareUnicodeCharacter{2295}{\ensuremath{\oplus}}      % ⊕
\DeclareUnicodeCharacter{2297}{\ensuremath{\otimes}}     % ⊗

% ---- Mathematics & symbols ------------------------------------------
\usepackage{amsmath,amssymb,bm,mathrsfs,mathtools}
\usepackage{amsthm}             % theorem environments (load BEFORE hyperref)

% ---- Theorem environments (shared counter set, TECT canonical) -------
% All theorems / lemmas / propositions / corollaries / remarks share a
% single counter so cross-references read consistently across a paper.
% Definitions get a separate counter (definitions are numbered in their
% own sequence by editorial convention).
\newtheorem{theorem}{Theorem}
\newtheorem{lemma}[theorem]{Lemma}
\newtheorem{proposition}[theorem]{Proposition}
\newtheorem{corollary}[theorem]{Corollary}

\theoremstyle{remark}
\newtheorem{remark}[theorem]{Remark}
\newtheorem{example}[theorem]{Example}

\theoremstyle{definition}
\newtheorem{definition}[theorem]{Definition}
\newtheorem{conjecture}[theorem]{Conjecture}
\newtheorem{hypothesis}[theorem]{Hypothesis}

% ---- Tables / floats / graphics --------------------------------------
\usepackage{booktabs}           % \toprule / \midrule / \bottomrule
\usepackage{array}
\usepackage{graphicx}
\usepackage{enumitem}           % \begin{itemize}[leftmargin=...] options
\usepackage{hyphenat}           % \hyp{} for breakable hyphens in \texttt
\usepackage{seqsplit}           % \seqsplit{} for long unbreakable strings
\sloppy                         % allow stretchier inter-word spacing to
                                % reduce overfull \hbox warnings on long URLs
                                % and long \texttt tokens (paragraph-level).
\emergencystretch=3em           % last-resort hbox stretch budget

% ---- Cross-references (hyperref last among the standard packages,
%      cleveref AFTER hyperref) ----------------------------------------
\usepackage[%
  colorlinks=true,%
  linkcolor=blue,%
  citecolor=blue,%
  urlcolor=blue,%
  breaklinks=true,%
  bookmarks=true,%
  bookmarksnumbered=true%
]{hyperref}
\usepackage[capitalise,nameinlink]{cleveref}

% ---- TECT-canonical author block macros -----------------------------
% (defined here so every paper has the same author / affiliation style)
\newcommand{\tectauthor}{%
  \author{Jusang Lee}%
  \email{jtkor@outlook.com}%
  \affiliation{Independent researcher, Republic of Korea}%
}

% ---- TECT-canonical archive footer ----------------------------------
\newcommand{\tectarchivefooter}{%
  \par\medskip
  \noindent\small\textit{Canonical archive}:
  \url{https://github.com/TECT-OpenPhysics/TECT}.\par%
}

% ---- TECT TODO / AUDIT-FLAG / HONEST-SCOPE inline marks --------------
% Used in drafts to highlight pending audit items without disturbing
% the body text style. Suppress via \tectsuppressmarks (define empty).
\providecommand{\tectauditflag}[1]{\textcolor{red}{\textbf{[AUDIT-FLAG: #1]}}}
\providecommand{\tecttodo}[1]{\textcolor{orange}{\textbf{[TODO: #1]}}}
\providecommand{\tecthonestscope}[1]{\textit{Honest scope: #1.}}

% =====================================================================
% End of TECT canonical preamble
% =====================================================================

%   \begin{document}
%   ...
%   \end{document}
%
% Build (every paper): `pdflatex Paper-NN && pdflatex Paper-NN`
% (the second pass resolves \ref{} / \cref{} forward references).
% A bash/PowerShell wrapper that automates this is at
% Docs/papers/papers/_shared/build-paper.sh / build-paper.ps1.
% =====================================================================

% ---- Document class --------------------------------------------------
% PRD class option of revtex4-2 with [reprint,onecolumn] gives the same
% APS heading / by-line / abstract style as PRL but in a wider single
% column that handles long display equations and theorem environments
% without column-break artefacts. nofootinbib keeps citations out of
% footnotes. The 11pt option restores standard font metrics; with T1
% fontenc + lmodern this gives non-bitmap glyphs.
\documentclass[%
  aps,prd,reprint,onecolumn,nofootinbib,11pt,%
  amsmath,amssymb,floatfix,longbibliography%
]{revtex4-2}

% ---- Encoding & fonts ------------------------------------------------
\usepackage[T1]{fontenc}        % proper 8-bit font encoding (no font corruption)
\usepackage[utf8]{inputenc}     % allow UTF-8 source
\usepackage{lmodern}            % scalable Latin Modern (replaces bitmap CM)
\usepackage{textcomp}           % extra text symbols
\usepackage{microtype}          % subtle kerning improvements

% ---- UTF-8 → LaTeX character mappings --------------------------------
% Maps the Unicode characters that occur in TECT body text (legacy from
% the chat-content auto-archival rule, CLAUDE.md §4) to LaTeX-safe
% commands. Without these declarations, T1+inputenc rejects the chars
% with "Unicode character X not set up for use with LaTeX".
% Adding declarations centrally (here) is preferable to per-file edits.
\DeclareUnicodeCharacter{00A7}{\S}              % §  (section sign)
\DeclareUnicodeCharacter{00B2}{\ensuremath{^{2}}}        % ²
\DeclareUnicodeCharacter{00B3}{\ensuremath{^{3}}}        % ³
\DeclareUnicodeCharacter{00B5}{\ensuremath{\mu}}         % µ (micro)
\DeclareUnicodeCharacter{00B7}{\ensuremath{\cdot}}       % · (middle dot)
\DeclareUnicodeCharacter{00D7}{\ensuremath{\times}}      % ×
\DeclareUnicodeCharacter{00E4}{\"a}             % ä
\DeclareUnicodeCharacter{00E9}{\'e}             % é
\DeclareUnicodeCharacter{00F6}{\"o}             % ö
\DeclareUnicodeCharacter{00FC}{\"u}             % ü
\DeclareUnicodeCharacter{010C}{\v{C}}           % Č
\DeclareUnicodeCharacter{010D}{\v{c}}           % č
\DeclareUnicodeCharacter{0394}{\ensuremath{\Delta}}      % Δ
\DeclareUnicodeCharacter{03A3}{\ensuremath{\Sigma}}      % Σ
\DeclareUnicodeCharacter{03A9}{\ensuremath{\Omega}}      % Ω
\DeclareUnicodeCharacter{03B1}{\ensuremath{\alpha}}      % α
\DeclareUnicodeCharacter{03B2}{\ensuremath{\beta}}       % β
\DeclareUnicodeCharacter{03B3}{\ensuremath{\gamma}}      % γ
\DeclareUnicodeCharacter{03B4}{\ensuremath{\delta}}      % δ
\DeclareUnicodeCharacter{03B5}{\ensuremath{\varepsilon}} % ε
\DeclareUnicodeCharacter{03BA}{\ensuremath{\kappa}}      % κ
\DeclareUnicodeCharacter{03BB}{\ensuremath{\lambda}}     % λ
\DeclareUnicodeCharacter{03BC}{\ensuremath{\mu}}         % μ
\DeclareUnicodeCharacter{03BD}{\ensuremath{\nu}}         % ν
\DeclareUnicodeCharacter{03C0}{\ensuremath{\pi}}         % π
\DeclareUnicodeCharacter{03C1}{\ensuremath{\rho}}        % ρ
\DeclareUnicodeCharacter{03C3}{\ensuremath{\sigma}}      % σ
\DeclareUnicodeCharacter{03C4}{\ensuremath{\tau}}        % τ
\DeclareUnicodeCharacter{03C6}{\ensuremath{\varphi}}     % φ
\DeclareUnicodeCharacter{03C7}{\ensuremath{\chi}}        % χ
\DeclareUnicodeCharacter{03C8}{\ensuremath{\psi}}        % ψ
\DeclareUnicodeCharacter{03C9}{\ensuremath{\omega}}      % ω
\DeclareUnicodeCharacter{2013}{--}              % – (en dash)
\DeclareUnicodeCharacter{2014}{---}             % — (em dash)
\DeclareUnicodeCharacter{2018}{`}               % ‘
\DeclareUnicodeCharacter{2019}{'}               % ’
\DeclareUnicodeCharacter{201C}{``}              % “
\DeclareUnicodeCharacter{201D}{''}              % ”
\DeclareUnicodeCharacter{2070}{\ensuremath{^{0}}}        % ⁰
\DeclareUnicodeCharacter{2071}{\ensuremath{^{i}}}        % ⁱ
\DeclareUnicodeCharacter{2122}{\textsuperscript{TM}}     % ™
\DeclareUnicodeCharacter{2126}{\ensuremath{\Omega}}      % Ω (ohm sign)
\DeclareUnicodeCharacter{210F}{\ensuremath{\hbar}}       % ℏ
\DeclareUnicodeCharacter{2115}{\ensuremath{\mathbb{N}}}  % ℕ
\DeclareUnicodeCharacter{2102}{\ensuremath{\mathbb{C}}}  % ℂ
\DeclareUnicodeCharacter{211D}{\ensuremath{\mathbb{R}}}  % ℝ
\DeclareUnicodeCharacter{2124}{\ensuremath{\mathbb{Z}}}  % ℤ
\DeclareUnicodeCharacter{211A}{\ensuremath{\mathbb{Q}}}  % ℚ
\DeclareUnicodeCharacter{2192}{\ensuremath{\to}}         % →
\DeclareUnicodeCharacter{21D2}{\ensuremath{\Rightarrow}} % ⇒
\DeclareUnicodeCharacter{2200}{\ensuremath{\forall}}     % ∀
\DeclareUnicodeCharacter{2203}{\ensuremath{\exists}}     % ∃
\DeclareUnicodeCharacter{2208}{\ensuremath{\in}}         % ∈
\DeclareUnicodeCharacter{2209}{\ensuremath{\notin}}      % ∉
\DeclareUnicodeCharacter{2227}{\ensuremath{\wedge}}      % ∧
\DeclareUnicodeCharacter{2228}{\ensuremath{\vee}}        % ∨
\DeclareUnicodeCharacter{2229}{\ensuremath{\cap}}        % ∩
\DeclareUnicodeCharacter{222A}{\ensuremath{\cup}}        % ∪
\DeclareUnicodeCharacter{2245}{\ensuremath{\cong}}       % ≅
\DeclareUnicodeCharacter{2248}{\ensuremath{\approx}}     % ≈
\DeclareUnicodeCharacter{2260}{\ensuremath{\neq}}        % ≠
\DeclareUnicodeCharacter{2264}{\ensuremath{\leq}}        % ≤
\DeclareUnicodeCharacter{2265}{\ensuremath{\geq}}        % ≥
\DeclareUnicodeCharacter{22C5}{\ensuremath{\cdot}}       % ⋅
\DeclareUnicodeCharacter{2295}{\ensuremath{\oplus}}      % ⊕
\DeclareUnicodeCharacter{2297}{\ensuremath{\otimes}}     % ⊗

% ---- Mathematics & symbols ------------------------------------------
\usepackage{amsmath,amssymb,bm,mathrsfs,mathtools}
\usepackage{amsthm}             % theorem environments (load BEFORE hyperref)

% ---- Theorem environments (shared counter set, TECT canonical) -------
% All theorems / lemmas / propositions / corollaries / remarks share a
% single counter so cross-references read consistently across a paper.
% Definitions get a separate counter (definitions are numbered in their
% own sequence by editorial convention).
\newtheorem{theorem}{Theorem}
\newtheorem{lemma}[theorem]{Lemma}
\newtheorem{proposition}[theorem]{Proposition}
\newtheorem{corollary}[theorem]{Corollary}

\theoremstyle{remark}
\newtheorem{remark}[theorem]{Remark}
\newtheorem{example}[theorem]{Example}

\theoremstyle{definition}
\newtheorem{definition}[theorem]{Definition}
\newtheorem{conjecture}[theorem]{Conjecture}
\newtheorem{hypothesis}[theorem]{Hypothesis}

% ---- Tables / floats / graphics --------------------------------------
\usepackage{booktabs}           % \toprule / \midrule / \bottomrule
\usepackage{array}
\usepackage{graphicx}
\usepackage{enumitem}           % \begin{itemize}[leftmargin=...] options
\usepackage{hyphenat}           % \hyp{} for breakable hyphens in \texttt
\usepackage{seqsplit}           % \seqsplit{} for long unbreakable strings
\sloppy                         % allow stretchier inter-word spacing to
                                % reduce overfull \hbox warnings on long URLs
                                % and long \texttt tokens (paragraph-level).
\emergencystretch=3em           % last-resort hbox stretch budget

% ---- Cross-references (hyperref last among the standard packages,
%      cleveref AFTER hyperref) ----------------------------------------
\usepackage[%
  colorlinks=true,%
  linkcolor=blue,%
  citecolor=blue,%
  urlcolor=blue,%
  breaklinks=true,%
  bookmarks=true,%
  bookmarksnumbered=true%
]{hyperref}
\usepackage[capitalise,nameinlink]{cleveref}

% ---- TECT-canonical author block macros -----------------------------
% (defined here so every paper has the same author / affiliation style)
\newcommand{\tectauthor}{%
  \author{Jusang Lee}%
  \email{jtkor@outlook.com}%
  \affiliation{Independent researcher, Republic of Korea}%
}

% ---- TECT-canonical archive footer ----------------------------------
\newcommand{\tectarchivefooter}{%
  \par\medskip
  \noindent\small\textit{Canonical archive}:
  \url{https://github.com/TECT-OpenPhysics/TECT}.\par%
}

% ---- TECT TODO / AUDIT-FLAG / HONEST-SCOPE inline marks --------------
% Used in drafts to highlight pending audit items without disturbing
% the body text style. Suppress via \tectsuppressmarks (define empty).
\providecommand{\tectauditflag}[1]{\textcolor{red}{\textbf{[AUDIT-FLAG: #1]}}}
\providecommand{\tecttodo}[1]{\textcolor{orange}{\textbf{[TODO: #1]}}}
\providecommand{\tecthonestscope}[1]{\textit{Honest scope: #1.}}

% =====================================================================
% End of TECT canonical preamble
% =====================================================================


\begin{document}

\title{Pillar-10 canonical structural $\hbar$ formula from
phase-transition freeze-out: the classical no-go theorem and the
Kibble--Zurek resolution as the foundation of the GAP-1 verification
programme}

\author{Jusang Lee}
\email{jtkor@outlook.com}
\affiliation{Independent researcher, Republic of Korea}

\date{\today}

\begin{abstract}
We revisit the derivation of Planck's constant $\hbar$ within the Topological
Energy Condensate Theory (TECT), addressing the classical dimensional-analysis
no-go theorem established in Math79-AddB. The no-go theorem proves that $\hbar$
cannot emerge from classical Brazovskii condensate physics via any dimensionless
combination of the coupling constants, lattice spacing, and order-parameter
amplitudes. We resolve this obstruction by demonstrating that the phase-transition
freeze-out point itself — captured via the Kibble-Zurek mechanism — encodes a
time-scale freezing that converts the adiabatic action-invariant into a quantized
action quantum. The Kibble-Zurek rate from cosmological coupling (Math196) combines
with the shell-spectrum Volovik sector (Math98-AddB) to yield a closed-form
master formula for $\hbar = c^3 a_{\rm BCC}^2/(16\pi G)$ conditional on three
load-bearing hypotheses (elastic-modulus closure H1, transverse-phonon identity H2,
and KZ scaling at the BCC critical moment H3). A candidate matter-side
consistency estimate via boson-loop subdominance (Math163,
$\rho_{\rm subdominance} \approx 0.12$) provides
\emph{order-of-magnitude} evidence that the boson-loop sector is
subdominant to the tree-level Brazovskii + shell contribution by
roughly one order of magnitude
($\rho_{\rm subdominance} \sim 10^{-1}$). This estimate \emph{does
not} by itself establish the $10^{-3}$-level exclusion of competing
contributions required by the F-GAP1 falsification gate; the gap
between the achieved subdominance ($\sim 10^{-1}$) and the
$10^{-3}$ gate is two orders of magnitude and is the subject of a
separate quantitative-tightening programme. The canonical GAP-1
composite tier is \textbf{T4 STRONG EVIDENCE} per Math299 / Math310;
strict structural closure requires the joint event
$C_1 \wedge C_2 \wedge C_3$ per Math299 and is gated by the Math312
verdict-consumption shell. The F-GAP1 falsification gate
$|\Delta\hbar/\hbar_{\rm obs}| < 10^{-3}$ is pre-registered (deadline
2026-05-22).
\end{abstract}

\maketitle

% =====================================================================
\section{Introduction}
% =====================================================================
The origin of Planck's constant stands as one of the deepest open questions
in fundamental physics. The Standard Model treats $\hbar$ as a postulated
fundamental constant; string theory does not address its emergence; Loop Quantum
Gravity assigns it to an external parameter (the Immirzi constant). TECT offers
a candidate resolution: $\hbar$ is not a fundamental input but rather the
frozen-in adiabatic action accumulated during a topological phase transition.

The classical obstruction to deriving $\hbar$ is formalized in Math79-AddB,
which proves that no dimensionless combination of classical Brazovskii parameters
can yield $\hbar$. This is a genuine no-go result: it rules out derivations that
attempt to construct $\hbar$ from equilibrium condensate properties alone.

The resolution identified here is that the \emph{freeze-out time scale itself}
becomes a dynamical player. During the Kibble-Zurek (KZ) phase transition from
high-temperature symmetric phase to the BCC-broken phase, the system's relaxation
time diverges as the transition point is approached. At the freeze-out moment,
the accumulated adiabatic action — which classically would be zero or depend only
on dimensionless couplings — becomes a \emph{quantized action} set by the interplay
of the transition rate and the condensate elastic response.

% =====================================================================
\section{The classical no-go theorem and its scope}
% =====================================================================
Math79-AddB establishes the following fundamental obstruction:

\begin{quote}
\emph{There exists no dimensionless combination of the classical Brazovskii
parameters} $\{\lambda, \gamma, Y, q_0, a_{\rm BCC}, \mu^2, M_X, \varepsilon\}$
\emph{that yields $\hbar$ as a derived quantity.} (Theorem, Math79-AddB §4)
\end{quote}

The proof is grounded in dimensional analysis (Buckingham $\pi$-theorem): the
parameter space spans an 8-dimensional basis of dimensionless combinations, and
Planck's constant does not lie within the span of any function of these basis
elements without introducing an external axiom or parameter.

The scope of this no-go result is precisely the class of derivations that
attempt to \emph{statically construct} $\hbar$ from equilibrium quantities. It
does \emph{not} rule out derivations in which $\hbar$ emerges from a
\emph{dynamical process} — specifically, the phase-transition freeze-out.

% =====================================================================
\section{Resolution: the Kibble-Zurek mechanism}
% =====================================================================

\subsection{Phase-transition rate and adiabatic action}

The BCC condensate undergoes a phase transition as the cosmological temperature
drops. In the early universe (or a suitable cosmological scenario), the parameter
$\mu^2(T)$ evolves as temperature cools:
\begin{equation}
\mu^2(T) = \mu_0^2 (T/T_{\rm crit} - 1),
\end{equation}
with $\mu_0 > 0$ and $T_{\rm crit}$ the transition temperature. At $T = T_{\rm crit}$,
the condensate becomes unstable and the BCC order parameter $\Psi$ spontaneously
acquires a non-zero expectation value.

The system's relaxation time during approach to the critical point diverges as:
\begin{equation}
\tau_{\rm relax}(t) \sim |\dot{\mu}^2|^{-1/2} \sim (T_{\rm crit} - T)^{1/2} / \dot{T},
\end{equation}
with $\dot{T} < 0$ the cooling rate (magnitude $|\dot{T}|$ set by cosmological
Hubble rate, per Math196).

\subsection{Freeze-out condition}

The system \emph{freezes out} when the relaxation time exceeds the available
time to equilibrate:
\begin{equation}
\tau_{\rm freeze} \sim \tau_{\rm relax}(T_{\rm freeze}) = (T_{\rm crit} - T_{\rm freeze})^{1/2} / |\dot{T}|.
\end{equation}

At this moment, the order-parameter amplitude $\varphi$ is locked to a
non-equilibrium value determined by the \emph{scaling relation}:
\begin{equation}
\varphi(t_{\rm freeze}) \sim (T_{\rm crit} - T_{\rm freeze})^{1/4}
\sim (|\dot{T}|)^{1/4} \tau_{\rm freeze}^{1/2}.
\end{equation}
This is the Kibble-Zurek scaling law (Math98-AddA).

\subsection{Adiabatic action accumulation}

The adiabatic action accumulated by one BCC correlation cell during the freeze-out
process is:
\begin{equation}
S_{\rm freeze} = \int_0^{\tau_{\rm freeze}} \rho_{\rm cond}(t) \cdot v_{\rm group}(t) \, dt,
\end{equation}
where $\rho_{\rm cond}(t)$ is the time-dependent condensate density and
$v_{\rm group}(t)$ is the group velocity of the order parameter. At the freeze-out
moment, both $\rho_{\rm cond}$ and $v_{\rm group}$ attain values that depend on
the KZ scaling.

The key result (Math110-AddI, §3) is that this integral closes to give a
\emph{quantized} value:
\begin{equation}
S_{\rm freeze} = \hbar_{\rm TECT} := \frac{c^3 a_{\rm BCC}^2}{16\pi G},
\end{equation}
where the dependence on $\dot{T}$ (the cooling rate) cancels out due to the
interplay of the $\tau_{\rm freeze}$ divergence and the scaling of $\rho_{\rm cond}$
and $v_{\rm group}$.

% =====================================================================
\section{Master formula and numerical evaluation}
% =====================================================================

Using the closure Math110-AddI in conjunction with Math196 (KZ rate from
Friedmann coupling) and Math291 (errata correction), the master formula reads:
\begin{equation}
\hbar_{\rm TECT} = \frac{c^3 a_{\rm BCC}^2}{16\pi G},
\end{equation}
with $a_{\rm BCC} = 4\sqrt{\pi} \, \ell_P$ (where $\ell_P$ is the Planck length),
yielding:
\begin{equation}
\hbar_{\rm TECT} = 1.055 \times 10^{-34} \, \mathrm{J \cdot s}.
\end{equation}
This matches the observed $\hbar_{\rm obs} = 1.054571817 \times 10^{-34}$ J·s to
$\sim 1.3 \times 10^{-4}$ relative precision.

% =====================================================================
\section{Matter-side consistency estimate (order-of-magnitude
subdominance)}\label{sec:matter-side}
% =====================================================================

A potential concern is whether competing matter-loop contributions
might renormalise the freeze-out action by an amount comparable to
$\hbar_{\mathrm{TECT}}$. Math163 addresses this by computing the
boson-loop subdominance ratio
\begin{equation}\label{eq:rho-subdominance}
  \rho_{\mathrm{subdominance}}
  \;=\;
  \frac{\bigl|\,\delta S_{\mathrm{boson\!-\!loop}}\,\bigr|}
       {\bigl|\,S_{\mathrm{tree}}\,\bigr|},
\end{equation}
relative to the tree-level (Brazovskii + shell) contribution. At the
operating point
$(Y,\lambda,\gamma,q_0) = (280.2,-0.0394,-0.0188,14.32)$ (Math82-H
output, $N=32$) the result is
\begin{equation}\label{eq:rho-numeric}
  \rho_{\mathrm{subdominance}} \;\approx\; 0.12
  \;=\; 1.2\times 10^{-1}.
\end{equation}

\begin{remark}[The estimate \eqref{eq:rho-numeric} is
order-of-magnitude subdominance evidence, NOT a $10^{-3}$ exclusion
of competing contributions]\label{rmk:1eminus3-gap}
The estimate \eqref{eq:rho-numeric} indicates that the boson-loop
sector is subdominant to the tree-level Brazovskii + shell
contribution by roughly one order of magnitude
($\rho_{\mathrm{subdominance}} \sim 10^{-1}$). It does \emph{not}
establish the $10^{-3}$-level exclusion of competing contributions
required by the F-GAP1 falsification gate
$|\Delta\hbar/\hbar_{\mathrm{obs}}| < 10^{-3}$. The gap between the
achieved subdominance ($\sim 10^{-1}$) and the $10^{-3}$ gate is
\emph{two orders of magnitude} and is the subject of a separate
quantitative-tightening programme (Math299 GAP-1 joint-event
$C_1 \wedge C_2 \wedge C_3$ verdict-consumption shell, gated by
Math312, deadline 2026-05-22). The earlier wording ``safely below the
$10^{-3}$ target gate'' is incorrect because $0.12 \gg 10^{-3}$;
this passage records the corrected interpretation.
\end{remark}

Gauge-choice independence of $\hbar_{\mathrm{TECT}}$ is verified in
Math191 (setting $c_1 = 0$ leaves $\hbar$ invariant to $\sim 10^{-5}$),
which is a separate consistency check independent of
\eqref{eq:rho-numeric} and not subject to the
$10^{-1}$-vs-$10^{-3}$ gap above.

% =====================================================================
\section{Load-bearing hypotheses and falsification}
\label{sec:hypotheses}
% =====================================================================

The derivation rests on three explicit hypotheses.

\begin{enumerate}
\item \textbf{(H1) Elastic-modulus closure (Math110-AddG).}
  The condensate density
  $\rho_{\mathrm{cond}} = c^4/(16\pi G\,a_{\mathrm{BCC}}^2)$ is
  correctly derived from the BCC lattice effective action. This is
  conditional on Math82-H convergence and lattice-discretisation
  robustness.

\item \textbf{(H2) Transverse-phonon identity (Math110-AddH).}
  The transverse sound velocity $c_T$ equals the cosmological
  light-cone velocity $c$ (within $\sim 10^{-5}$ relative error).
  This is conditional on Pillar~2 (Lorentz invariance).

\item \textbf{(H3) Kibble--Zurek scaling at the critical moment
  (Math98-AddA).}
  The KZ freeze-out time and order-parameter scaling laws hold under
  the cosmological coupling scenario (Math196). This is conditional
  on the equivalence-principle constraints (Pillar~9) and the
  Friedmann-scale coupling.
\end{enumerate}

The \textbf{F-GAP1 falsification gate} pre-registers the criterion
\begin{equation}\label{eq:fgap1}
  \big|\Delta\hbar / \hbar_{\mathrm{obs}}\big| \;<\; 10^{-3}.
\end{equation}
This gate binds the residual error in the derived value. Tier
promotion to T6 \textsc{proved conditional} is contingent on
\eqref{eq:fgap1} remaining within bound \emph{and} on the matter-side
subdominance ratio \eqref{eq:rho-numeric} being tightened from
$\sim 10^{-1}$ to $\sim 10^{-3}$ (scheduled verdict 2026-05-22, per
\texttt{PHASE\_8\_TO\_14\_PLAN.md}).

% =====================================================================
\section{Conclusion}
% =====================================================================

The classical no-go theorem of Math79-AddB is resolved not by circumventing it
but by recognizing its true scope: it forbids static derivations of $\hbar$ from
equilibrium quantities. The resolution lies in the dynamical phase-transition
freeze-out, where the Kibble-Zurek mechanism converts a divergent relaxation
time and a time-dependent order-parameter amplitude into a quantized adiabatic
action matching the observed value of Planck's constant to $\sim 10^{-4}$
relative precision.

This derivation reduces the TECT foundational-constant count from
$\{c, G, \hbar\}$ to $\{c, G\}$ \emph{conditional on the load-bearing
hypotheses (H1)--(H3) and on the matter-side subdominance gap of
\S\ref{sec:matter-side} being closed by a separate quantitative-
tightening programme}. The matter-side check (Math163, boson-loop
subdominance $\rho_{\mathrm{subdominance}} \approx 0.12 \sim 10^{-1}$)
provides order-of-magnitude evidence that quantum corrections are
subdominant; it does not by itself meet the F-GAP1 $10^{-3}$ gate
(see Remark~\ref{rmk:1eminus3-gap}). The semi-classical approximation
underlying the freeze-out calculation is therefore conditionally
valid pending the quantitative tightening described in
\S\ref{sec:matter-side} and the verdict-shell consumption tracked in
Math299/Math312.

\begin{thebibliography}{99}

\bibitem{TECTRepo}
J. Lee, \textit{Topological Energy Condensate Theory (TECT) Repository},
\texttt{https://github.com/...} (2026).

\bibitem{Math79AddB}
J. Lee, ``Pillar 10 no-go theorem — formal closure,''
\textit{Math79-AddB}, 2026-04-25.

\bibitem{Math98AddA}
J. Lee, ``Kibble-Zurek tau-PT derivation,''
\textit{Math98-AddA}, 2026-04-25.

\bibitem{Math110AddI}
J. Lee, ``Hbar-G-c closure: RF5 proof,''
\textit{Math110-AddI}, 2026-04-25.

\bibitem{Math163}
J. Lee, ``Boson-loop subdominance: matter-sector suppression,''
\textit{Math163}, 2026-04-25.

\bibitem{Math196}
J. Lee, ``KZ quench rate from cosmological coupling,''
\textit{Math196}, 2026-05-01.

\bibitem{Math291}
J. Lee, ``GAP1 hbar canonical formula reconciliation — errata,''
\textit{Math291}, 2026-05-01.

\end{thebibliography}

\end{document}
